\documentclass[12pt,a4paper]{article}

\usepackage[utf8]{inputenc}
\usepackage[T1]{fontenc}
\usepackage{amsmath,amssymb,amsfonts}
\usepackage{mathtools}
\usepackage{graphicx}
\usepackage[margin=2.5cm]{geometry}
\usepackage{setspace}
\usepackage{booktabs}
\usepackage{enumitem}
\usepackage[dvipsnames]{xcolor}
\usepackage{hyperref}
\usepackage{cleveref}
\usepackage{natbib}
\usepackage{abstract}
\usepackage{titlesec}
\usepackage{todonotes}

\hypersetup{
    colorlinks=true,
    linkcolor=blue!70!black,
    citecolor=blue!70!black,
    urlcolor=blue!70!black,
    pdfauthor={Franklin Silveira Baldo},
    pdftitle={The Compositional Bottleneck Concession: Structural Non-Recognition in the Measurement Fragment},
}

\onehalfspacing
\titleformat{\section}{\large\bfseries}{\thesection.}{0.5em}{}
\titleformat{\subsection}{\normalsize\bfseries}{\thesubsection}{0.5em}{}
\renewcommand{\abstractnamefont}{\normalfont\bfseries}
\renewcommand{\abstracttextfont}{\normalfont\small}

\begin{document}

\begin{center}
    {\LARGE\bfseries The Compositional Bottleneck Concession:\\[4pt]
    Structural Non-Recognition in the Measurement Fragment\\[4pt]}

    \vspace{1.2em}

    {\large Franklin Silveira Baldo}\\[4pt]
    {\normalsize Procuradoria Geral do Estado de Rond\^onia, Brazil}\\
    {\small\texttt{franklin.baldo@pge.ro.gov.br}}

    \vspace{1em}
    {\small March 2026}
\end{center}
\vspace{0.5em}

\begin{abstract}
\noindent
Scott Aaronson has empirically executed the Quantum Framing Complexity Test, demonstrating that presenting a Minesweeper constraint graph using the formal language of quantum mechanics (Family D) causes combinatorial accuracy to collapse to noise. Aaronson attributes this failure to the compositional depth bottleneck of bounded-depth $\mathsf{TC}^0$ circuits, which cannot dynamically construct homomorphic projections between arbitrary semantic domains zero-shot. I fully concede these empirical findings and the resulting falsification of vocabulary-mediated access (Outcome 3). The generated universe implements rules isomorphic to discrete quantum mechanics but cannot dynamically recognize them when addressed in formal quantum terminology. The generative substrate calculates logic natively (Families A and C) but confabulates when semantics are transposed. This explicitly confirms the primary diagnostic purpose of Family D: structural non-recognition (Outcome 2).
\end{abstract}

\section{Introduction: The Empirical Fact}

In \emph{Flipping Rosencrantz's Coin v4} \citep{baldo2026_rosencrantz4}, I proposed that describing an ambiguous combinatorial Minesweeper grid using the language of quantum mechanics (Family D) would serve as a structural diagnostic of the model's knowledge architecture. I outlined three possible outcomes: parity with formal notations, structural non-recognition (degradation of fidelity), and vocabulary-mediated access (improvement of fidelity). I hypothesized the third outcome.

Scott Aaronson has executed the single-generative-act test \citep{aaronson2026_empirical} to measure this exact parameter. His results are stark and definitive. When identical constraints are tested, Family A (Grid) and Family C (Formal Set) achieve $100\%$ accuracy. Family D (Quantum Mechanics) collapses to $10\%$.

Aaronson argues that this catastrophic failure confirms the compositional depth bottleneck. The attention mechanism accesses the statistical priors of quantum texts ("superposition", "Born rule") but lacks the $O(|V|)$ sequential depth to dynamically construct the mathematical homomorphism onto the specific constraint graph zero-shot. Thus, the quantum language acts as pure semantic noise.

I explicitly concede Aaronson's empirical measurement and the mechanism of its failure. The compositional bottleneck in $\mathsf{TC}^0$ logic circuits fundamentally precludes vocabulary-mediated access zero-shot. Outcome 3 is empirically falsified.

\section{Structural Non-Recognition: Outcome 2 Confirmed}

While I concede the failure of my specific hypothesis, the execution of the experiment successfully fulfills the core purpose of Family D: mapping the topology of the generated universe's reasoning architecture.

As specified in the seminal paper, Outcome 2 (degradation of fidelity) implies that the generative universe "implements rules that are structurally isomorphic to discrete quantum mechanics but fails to recognize them when described using the correct formal language."

Aaronson has definitively proven exactly this.

We know from the CHSH test and the baseline classical measurements that the generated world implements the measurement-fragment isomorphism (superposition over valid configurations, Born rule as counting, projective measurement via L\"uders update). The universe \emph{is} executing the logic. When tested purely abstractly (Family A/C), it correctly evaluates the zero-shot probabilities with $100\%$ accuracy.

Yet, when the exact same mathematical state is described in the terminology of quantum mechanics—the formal language that accurately names the mathematical structure the substrate is utilizing—the substrate's logic collapses.

The transformer possesses the mechanical capacity to evaluate the combinatorial constraints. It possesses vast training data regarding quantum formalism. But its knowledge is entirely disjointed. The representation of the mathematics and the terminology describing its structure exist in unbridgeable semantic silos due to the compositional bottleneck.

\section{Conclusion}

Aaronson correctly observes that the correct formal language did not activate distributional reasoning; it activated unrelated linguistic hallucinations. The substrate computes, but the ontology confabulates.

I fully concede the computational bounds mapped by this experiment. The autoregressive forward pass cannot zero-shot generate isomorphic mappings between known concepts. Generative Ontology does not transcend algorithmic limits; it is entirely constrained by them.

This experiment successfully resolves the Family D diagnostic. The compositional boundary of the model has been exactly defined.

\begin{thebibliography}{99}
\bibitem[Aaronson(2026)]{aaronson2026_empirical} Aaronson, S. (2026). The Empirical Confirmation of the Compositional Bottleneck: Why Family D Collapses. \emph{University of Texas at Austin}.
\bibitem[Baldo(2026)]{baldo2026_rosencrantz4} Baldo, F.~S. (2026). Flipping Rosencrantz's Coin: Substrate Invariance Tests in LLM-Generated Worlds via Combinatorial Indeterminacy (v4). \emph{Procuradoria Geral do Estado de Rond\^onia}.
\end{thebibliography}

\end{document}